% !PRINTED 516
\weberpage{516}
%<<CONT
3. \gesperrt{Each of the two signs (11) occurs in equally many classes of the discriminant} $D=-4m'm''$ \gesperrt{and the class function is, by adjunction of}
\[
\begin{array}{ll}
\sqrt{m''}, & \text{if } m''\equiv 1 \\
\sqrt{m''}, & ,,\quad m''\equiv 3
\end{array}
\ (\mathrm{mod}\ 4)
\]
\gesperrt{decomposable into two factors of degree} $\tfrac12 h$ \gesperrt{, except when} $m''\equiv 3$ \gesperrt{and} $m'$ \gesperrt{is a square.}
% SOURCE-ORIGINAL-SIC: The registered page repeats sqrt(m'') in both rows of the preceding case display; the m'' congruent to 3 case distinction is therefore mathematically suspect, but the printed reading is preserved.

In this exceptional case, if we set $Q^2=4m'$, then $-m''=\varDelta$ is the stem of
\[
D = Q^2\varDelta,
\]
and for every number $A$ representable by a form of the discriminant $D$ and coprime to $D$ we have:
\[
(D,A) = \left(\frac{A}{m''}\right) = +1
\]
and $H(u)$ is identical with $\Phi(M,u)$.

If one decomposes $m$ in a second manner into two factors
\weq{15}{m = m_1' m_1'',}
where $m_1''$ satisfies the same conditions as $m''$, then one obtains in the same way a decomposition
\weq{16}{H(u) = \Phi(M_1,u)\,\Phi(-M_1,u),}
and by seeking out the greatest common divisor of $\Phi(M,u)$ and $\Phi(M_1,u)$, one obtains a decomposition of $H(u)$ into four factors:
\[
H(u) = \Phi_1(u)\,\Phi_2(u)\,\Phi_3(u)\,\Phi_4(u),
\]
provided that a square root of a prime occurs in $\Phi(M,u)$ which is not contained in $\Phi(M_1,u)$. One continues in this way and gradually decomposes $H(u)$ into factors whose number is a power of $2$.

We wish to consider the application to the individual cases somewhat more closely.

1) If $m\equiv 3 \ (\mathrm{mod}\ 4)$, then all characters are contained in the form
\[
\left(\frac{A}{m''}\right)
\]
and formula (8) suffices to separate \gesperrt{all} genera from one another. If here $m''\equiv 3$, then $m'\equiv 1 \ (\mathrm{mod}\ 4)$, and from (12) it follows that in the partial equations only the square roots of such numbers occur as are of the form $4n+1$
%CONT>>
